\documentclass[11pt]{article}

\usepackage{fontspec}
\usepackage[spanish,es-noshorthands]{babel}
\usepackage[margin=2.5cm]{geometry}
\usepackage{xurl}
\usepackage{hyperref}
\usepackage{array}
\usepackage{graphicx}
\usepackage{listings}

\setmainfont{Latin Modern Roman}
\setlength{\emergencystretch}{2em}

\lstset{
  basicstyle=\ttfamily\small,
  columns=fullflexible,
  keepspaces=true,
  breaklines=true,
  frame=single,
  framesep=4pt,
  xleftmargin=6pt,
  xrightmargin=6pt
}

% Captura de Postman: usa capturas/#1 si existe; si no, deja el recuadro para tomarla
\graphicspath{{capturas/}}
\newcommand{\captura}[2]{%
  \begin{center}
  \IfFileExists{capturas/#1}%
    {\includegraphics[width=\linewidth,height=0.4\textheight,keepaspectratio]{#1}}%
    {\fbox{\parbox[c][3cm][c]{0.9\linewidth}{\centering\small Captura pendiente: \texttt{capturas/#1}}}}\\[0.6em]
  {\small #2}
  \end{center}}

\hypersetup{
  pdftitle={Laboratorio 4. REST APIs},
  pdfauthor={Daniel Alonso Gracia Pinto},
  colorlinks=true,
  linkcolor=black,
  urlcolor=blue
}

\title{Laboratorio 4: REST APIs\\[0.35em]
\large HTTP, consumo desde Java, protocolos y diseño de una API con OpenAPI\\[0.25em]
\normalsize Ingeniería de Software II (2016702)}
\author{Daniel Alonso Gracia Pinto}
\date{Noviembre de 2026\\[0.35em]
\normalsize Grupo 00}

\begin{document}
\maketitle
\vspace{-1.6em}

\section*{Ejercicio 0. Anatomía de una petición HTTP}

Las peticiones van a \texttt{https://httpbin.org}, que devuelve en el cuerpo lo que recibió: los parámetros, los headers y el cuerpo de la petición.

\subsection*{(a) Partes de la URL}

La petición es \texttt{GET https://httpbin.org/get?curso=seii\&grupo=3}.

\begin{center}
\begin{tabular}{@{}p{4.2cm}p{5cm}@{}}
\hline
Parte & Valor \\
\hline
Esquema & \texttt{https} \\
Host & \texttt{httpbin.org} \\
Ruta & \texttt{/get} \\
Parámetros de consulta & \texttt{curso=seii}, \texttt{grupo=3} \\
\hline
\end{tabular}
\end{center}

El signo \texttt{?} separa la ruta de los parámetros y \texttt{\&} separa un parámetro del siguiente. En la respuesta, los parámetros aparecen en el campo \texttt{args}, ya convertidos en un objeto: \texttt{\{"curso": "seii", "grupo": "3"\}}. El \texttt{3} llega como texto, porque en la URL todo es texto y el servidor no sabe que era un número.

\captura{ej0-a.png}{Figura 1. \texttt{GET /get} con los parámetros \texttt{curso} y \texttt{grupo}.}

\subsection*{(b) Headers}

Además del \texttt{X-Equipo: Musubi Japan}, Postman agregó por su cuenta siete headers: \texttt{Cache-Control} (\texttt{no-cache}), \texttt{Postman-Token} (un identificador distinto en cada envío), \texttt{Host}, \texttt{User-Agent} (\texttt{PostmanRuntime/2.12.2}), \texttt{Accept} (\texttt{*/*}), \texttt{Accept-Encoding} (\texttt{gzip, deflate, br}) y \texttt{Connection} (\texttt{keep-alive}). El campo \texttt{headers} de la respuesta los muestra todos menos \texttt{Connection}: ese header solo vale entre dos equipos vecinos de la red, y el intermediario que tiene httpbin.org adelante no lo reenvía.

\begin{description}
  \raggedright
  \item[\texttt{User-Agent}.] Identifica el programa que hace la petición. El servidor lo usa para estadísticas, para adaptar la respuesta o para bloquear clientes. Con \texttt{curl} la misma petición llega con \texttt{curl/8.18.0}.
  \item[\texttt{Accept}.] Dice qué formatos de respuesta acepta el cliente. \texttt{*/*} significa cualquiera; una app que solo entiende JSON debería enviar \texttt{application/json}.
  \item[\texttt{Host}.] Dice a qué sitio va la petición. Un mismo servidor, con una sola IP, puede atender muchos dominios, y por el \texttt{Host} sabe cuál se pidió. En HTTP/1.1 es obligatorio.
\end{description}

\captura{ej0-b.png}{Figura 2. La misma petición con el header \texttt{X-Equipo}.}

\subsection*{(c) Cuerpo JSON}

El cuerpo enviado a \texttt{POST https://httpbin.org/post} es:

\begin{quote}
\ttfamily
\{\\
\hspace*{1em}"equipo": "Musubi Japan",\\
\hspace*{1em}"grupo": 0,\\
\hspace*{1em}"integrantes": ["Daniel Alonso Gracia Pinto"],\\
\hspace*{1em}"proyecto": "App móvil para aprender japonés"\\
\}
\end{quote}

El header que le dice al servidor que el cuerpo es JSON es \texttt{Content-Type: application/json}. Postman lo pone solo al elegir \emph{raw} y \emph{JSON}. La respuesta trae el cuerpo dos veces: en \texttt{data}, como texto, y en \texttt{json}, ya convertido en objeto.

Al cambiar a \emph{Text}, Postman envía \texttt{Content-Type: text/plain}. El httpbin.org de hoy sigue llenando \texttt{json} con ese header: intenta leer el cuerpo como JSON sin mirar el \texttt{Content-Type}. Lo comprobé con un cuerpo que no es JSON (\texttt{hola equipo}): ahí \texttt{json} sale en \texttt{null} y el texto solo queda en \texttt{data}. El httpbin original, en cambio, devolvía \texttt{json: null} siempre que el tipo no fuera \texttt{application/json}. Eso es lo que debería hacer una API real: no adivinar el formato, sino responder \texttt{415 Unsupported Media Type} si el \texttt{Content-Type} no es el que espera.

\captura{ej0-c-json.png}{Figura 3. \texttt{POST /post} con el cuerpo como JSON.}

\captura{ej0-c-text.png}{Figura 4. El mismo cuerpo enviado como \emph{Text}: \texttt{json} sigue lleno.}

\captura{ej0-c-text-hola.png}{Figura 5. Un cuerpo \emph{Text} que no es JSON: \texttt{json} sale en \texttt{null}.}

\subsection*{(d) Códigos de estado}

Las tres respuestas llegan con el cuerpo vacío (\texttt{Content-Length: 0}). La app solo cuenta con el código para decidir qué mostrar.

\begin{center}
\begin{tabular}{@{}>{\raggedright\arraybackslash}p{3.2cm}>{\raggedright\arraybackslash}p{1.2cm}>{\raggedright\arraybackslash}p{4.4cm}>{\raggedright\arraybackslash}p{5.3cm}@{}}
\hline
Código & Familia & Responsable & Qué muestra la app \\
\hline
201 Created & 2xx & No hay error: el servidor creó el recurso. & Una confirmación (``Préstamo registrado'') y el recurso nuevo. \\
404 Not Found & 4xx & El cliente: pidió un recurso que no existe o una URL mal armada. & ``No encontramos lo que buscas'' y la opción de volver. Reintentar no sirve. \\
500 Internal Server Error & 5xx & El servidor: falló al procesar una petición válida. & ``Algo salió mal de nuestro lado, intenta más tarde'', con un botón para reintentar y sin detalles técnicos. \\
429 Too Many Requests & 4xx & El cliente: superó el límite de peticiones del servidor. & Nada de inmediato: espera lo que indique el servidor y reintenta sola. Si tarda, ``Hay mucho tráfico, reintentando\ldots''. \\
\hline
\end{tabular}
\end{center}

El 429 no se pidió: apareció al enviar la petición de (a). httpbin.org admite 17 peticiones por segundo por IP. Postman estaba usando su \emph{Cloud Agent}, que envía la petición desde los servidores de Postman, con una IP que comparten muchos usuarios a la vez; la misma petición enviada con \texttt{curl} desde el equipo respondió 200. Al pasar al agente de escritorio, la petición salió con la IP propia y el 429 desapareció. A diferencia de los códigos de \texttt{/status}, este trae cuerpo, con el motivo y el tiempo de espera: \texttt{retry after 1 second}. Reintentar sí sirve aquí, pero esperando; reintentar enseguida solo alarga el bloqueo.

\captura{ej0-d-201.png}{Figura 6. Respuesta de \texttt{/status/201}.}

\captura{ej0-d-404.png}{Figura 7. Respuesta de \texttt{/status/404}.}

\captura{ej0-d-500.png}{Figura 8. Respuesta de \texttt{/status/500}.}

\captura{ej0-d-429.png}{Figura 9. El 429 que devolvió httpbin.org al superar el límite de peticiones.}

\section*{Ejercicio 1. CRUD completo con JSONPlaceholder}

El environment define \texttt{baseUrl = https://jsonplaceholder.typicode.com}, \texttt{postId = 1} y \texttt{userId = 1}. Las peticiones usan esas variables en las rutas y en los cuerpos; la tabla muestra las URLs ya resueltas. JSONPlaceholder acepta escrituras, pero no las guarda: responde lo que respondería un servidor real.

\subsection*{(a) Métodos, códigos y propiedades}

Un método es \emph{seguro} si no cambia nada en el servidor, e \emph{idempotente} si enviarlo una vez o varias deja el servidor en el mismo estado. Todo método seguro es idempotente; al revés no.

\begin{center}
\begin{tabular}{@{}>{\raggedright\arraybackslash}p{1.6cm}>{\raggedright\arraybackslash}p{3.3cm}>{\raggedright\arraybackslash}p{1.1cm}>{\raggedright\arraybackslash}p{1.3cm}>{\raggedright\arraybackslash}p{2cm}>{\raggedright\arraybackslash}p{4.4cm}@{}}
\hline
Método & Endpoint & Código & ¿Seguro? & ¿Idempotente? & Qué hace \\
\hline
GET & \texttt{/users} & 200 & Sí & Sí & Lista los 10 usuarios. \\
GET & \texttt{/users/1/posts} & 200 & Sí & Sí & Lista los 10 posts del usuario 1. \\
GET & \texttt{/posts?userId=1} & 200 & Sí & Sí & Filtra los posts por \texttt{userId}; devuelve los mismos 10. \\
POST & \texttt{/posts} & 201 & No & No & Crea un post. Responde con \texttt{id: 101} y \texttt{Location: /posts/101}. \\
PUT & \texttt{/posts/1} & 200 & No & Sí & Reemplaza el post 1 entero por el cuerpo enviado. \\
PATCH & \texttt{/posts/1} & 200 & No & No\textsuperscript{*} & Cambia solo los campos enviados y conserva los demás. \\
DELETE & \texttt{/posts/1} & 200 & No & Sí & Borra el post 1. Responde \texttt{\{\}}. \\
GET & \texttt{/posts/9999} & 404 & Sí & Sí & Busca un post que no existe. Responde \texttt{\{\}}. \\
\hline
\end{tabular}
\end{center}

{\small \textsuperscript{*}HTTP no garantiza que \texttt{PATCH} sea idempotente. Lo es cuando fija valores, como aquí (\texttt{title = "Solo cambia el titulo"}), y deja de serlo con operaciones relativas, como ``sumar 1 a las vistas''.}

\texttt{DELETE} es idempotente aunque la segunda vez responda distinto (en una API real, 404): lo que cuenta es el estado del servidor, y después de uno o de diez \texttt{DELETE} el post sigue sin existir. Una API real respondería \texttt{204 No Content} en lugar de \texttt{200} con un cuerpo vacío.

\captura{ej1-a-users.png}{Figura 10. \texttt{GET /users}.}

\captura{ej1-a-users-posts.png}{Figura 11. \texttt{GET /users/1/posts}.}

\captura{ej1-a-posts-userid.png}{Figura 12. \texttt{GET /posts?userId=1}.}

\captura{ej1-a-post.png}{Figura 13. \texttt{POST /posts}: el cuerpo de la respuesta.}

\captura{ej1-a-post-location.png}{Figura 14. \texttt{POST /posts}: el header \texttt{Location} de la respuesta.}

\captura{ej1-a-put.png}{Figura 15. \texttt{PUT /posts/1}.}

\captura{ej1-a-patch.png}{Figura 16. \texttt{PATCH /posts/1}.}

\captura{ej1-a-delete.png}{Figura 17. \texttt{DELETE /posts/1}.}

\captura{ej1-a-9999.png}{Figura 18. \texttt{GET /posts/9999}.}

\subsection*{(b) Ruta anidada o parámetro de consulta}

Las dos peticiones devuelven los mismos 10 posts, pero dicen cosas distintas. La ruta \texttt{/users/1/posts} dice que los posts \emph{pertenecen} al usuario 1: es una relación de pertenencia. El parámetro \texttt{/posts?userId=1} dice que de la colección de todos los posts se quieren solo algunos: es un filtro.

El criterio general es que la ruta identifica un recurso y la consulta filtra una colección. Conviene la ruta anidada cuando el hijo no tiene sentido sin el padre y casi siempre se consulta desde él, como las líneas de un pedido (\texttt{/pedidos/7/lineas}); conviene un nivel de anidamiento, dos como máximo. Conviene el parámetro cuando el criterio es opcional o se combina con otros: \url{/posts?userId=1&q=lab&page=2} se arma sin crear rutas nuevas, mientras que con rutas habría que inventar una por cada combinación.

\subsection*{(c) Crear y luego consultar}

El \texttt{POST /posts} respondió \texttt{201 Created}, un cuerpo con \texttt{id: 101} y el header \texttt{Location:} \url{https://jsonplaceholder.typicode.com/posts/101}. El \texttt{GET /posts/101} siguiente respondió \texttt{404 Not Found}.

JSONPlaceholder solo tiene 100 posts y no guarda las escrituras: calcula el \texttt{id} que tendría el post y arma la respuesta que daría un servidor real, pero no lo almacena. Por eso la URL del \texttt{Location} lleva a un recurso que no existe. Una prueba guarda el \texttt{id} en la variable \texttt{nuevoId}, y la petición siguiente consulta \texttt{\{\{baseUrl\}\}/posts/\{\{nuevoId\}\}}, sin escribir el 101 a mano.

En una API real, el \texttt{POST} debe responder \texttt{201 Created} con el header \texttt{Location} apuntando al recurso nuevo, y el \texttt{GET} a esa URL debe devolverlo con \texttt{200}. JSONPlaceholder imita bien la respuesta; lo que falta es el almacenamiento.

\captura{ej1-c-get-nuevo.png}{Figura 19. \texttt{GET /posts/\{\{nuevoId\}\}} después del \texttt{POST}: 404.}

\subsection*{(d) PUT y PATCH}

El \texttt{PUT /posts/1} envía el recurso completo:

\begin{lstlisting}
{ "id": 1, "title": "Titulo nuevo", "body": "Cuerpo nuevo", "userId": 1 }
\end{lstlisting}

El \texttt{PATCH /posts/1} envía solo lo que cambia:

\begin{lstlisting}
{ "title": "Solo cambia el titulo" }
\end{lstlisting}

La respuesta del \texttt{PATCH} trae el título nuevo y el \texttt{body} original del post 1 (\texttt{"quia et suscipit\ldots"}), porque los campos que no se enviaron se conservan.

\texttt{PUT} significa ``deja el recurso exactamente así''. Si se omite \texttt{title}, el post queda sin título: JSONPlaceholder lo muestra, porque su respuesta al \texttt{PUT} sin \texttt{title} trae solo \texttt{id}, \texttt{body} y \texttt{userId}. En una API real pasaría una de dos cosas: si \texttt{title} es obligatorio, el servidor rechaza la petición con \texttt{400 Bad Request}; si es opcional, lo borra o lo deja en \texttt{null}. En ningún caso conserva el título anterior: para cambiar un solo campo se usa \texttt{PATCH}.

\captura{ej1-d-put-sin-title.png}{Figura 20. \texttt{PUT /posts/1} sin \texttt{title}: la respuesta no lo trae.}

\subsection*{(e) Pruebas automáticas}

Las diez peticiones de la carpeta \texttt{Ejercicio 1} tienen pruebas en \emph{Scripts → Post-response}. Todas verifican el código de estado; además, cada una revisa lo propio de su método: que los filtros solo traigan posts del usuario 1, que el \texttt{POST} incluya \texttt{Location}, que el \texttt{PATCH} conserve el \texttt{body} y que el \texttt{PUT} sin \texttt{title} lo pierda. El archivo completo es \texttt{postman/scripts-ejercicio1.js}. Las de \texttt{GET /users} son estas:

\lstinputlisting[firstline=4,lastline=10]{../postman/scripts-ejercicio1.js}

El Collection Runner ejecutó la carpeta completa, en orden, con el environment \texttt{Lab4-00}: las 19 pruebas pasaron, sin errores, en 2,8 segundos. El orden importa: el \texttt{POST} va antes del \texttt{GET /posts/\{\{nuevoId\}\}} porque es el que define \texttt{nuevoId}.

\captura{ej1-e-runner.png}{Figura 21. Resultado del Collection Runner sobre la carpeta \texttt{Ejercicio 1}.}

\section*{Ejercicio 2. Un cliente REST en Java}

El proyecto Maven está en la carpeta \texttt{cliente-java} del repositorio. Hace las peticiones con \texttt{java.net.http.HttpClient} y convierte el JSON con Gson 2.11.0. El paquete \texttt{dominio} tiene los registros \texttt{Clima} y \texttt{Pais}, las interfaces de los repositorios y las excepciones; \texttt{infraestructura} tiene las clases que hablan HTTP. El dominio no importa nada de \texttt{java.net} ni de Gson.

\textbf{REST Countries ya no sirve.} La URL de la guía, \url{https://restcountries.com/v3.1/name/colombia?fields=name,capital,population,latlng}, hoy responde \texttt{301 Moved Permanently} hacia otro servidor, y ese servidor responde \texttt{200} con \texttt{"success": false} y el mensaje \emph{This API version has been deprecated}. Responde igual para Colombia y para Wakanda. La versión nueva, la v5, exige una llave: sin ella responde \texttt{401} con \texttt{authKeyMissing}, y con la llave de demostración devuelve siempre a Canadá. Por eso el programa usa \texttt{countries.dev}, un servicio gratuito y sin llave que conserva las rutas de REST Countries (\texttt{/name/\{nombre\}}, con \texttt{?fields=}). Lo que cambia es la forma del JSON: \texttt{name} y \texttt{capital} son texto en vez de objeto y lista. La clase se llama \texttt{CountriesDevPaisRepository}, y no \texttt{RestCountriesPaisRepository}, porque lleva el nombre de la API con la que habla.

\captura{ej2-restcountries-v31.png}{Figura 22. REST Countries v3.1 hoy: la versión está descontinuada.}

\subsection*{(a) Los repositorios}

En \texttt{OpenMeteoClimaRepository}, los cuatro TODO quedaron así:

\begin{lstlisting}
try {
    HttpResponse<String> res = http.send(req, HttpResponse.BodyHandlers.ofString());
    if (res.statusCode() != 200) {
        throw new ClimaNoDisponible("Open-Meteo respondió " + res.statusCode() + motivo(res.body()));
    }
    JsonObject actual = JsonParser.parseString(res.body())
        .getAsJsonObject().getAsJsonObject("current");
    return new Clima(actual.get("temperature_2m").getAsDouble(),
                     actual.get("wind_speed_10m").getAsDouble());
} catch (HttpTimeoutException e) {
    throw new ClimaNoDisponible("Open-Meteo no respondió en 5 segundos", e);
} catch (IOException e) {
    throw new ClimaNoDisponible("No hubo conexión con Open-Meteo", e);
} catch (InterruptedException e) {
    Thread.currentThread().interrupt();
    throw new ClimaNoDisponible("La consulta del clima se interrumpió", e);
} catch (JsonParseException | IllegalStateException
         | UnsupportedOperationException | NullPointerException e) {
    throw new ClimaNoDisponible("Open-Meteo respondió en un formato inesperado", e);
}
\end{lstlisting}

La conexión y la respuesta tienen cada una un límite de 5 segundos. \texttt{HttpTimeoutException} se atrapa antes que \texttt{IOException} porque es una subclase, y así el mensaje distingue ``tardó demasiado'' de ``no hay red''. \texttt{InterruptedException} vuelve a marcar el hilo como interrumpido antes de traducirse, para que quien llamó se entere. El último \texttt{catch} cubre un cuerpo que no es JSON o al que le falta \texttt{current}. Cuando Open-Meteo responde \texttt{400}, explica el motivo en el campo \texttt{reason}, y el mensaje lo incluye.

\texttt{PaisRepository} declara dos excepciones distintas: \texttt{PaisNoEncontrado} (el servicio respondió, pero el país no existe) y \texttt{PaisNoDisponible} (el servicio falló). La diferencia importa porque lo que hace el usuario es distinto en cada caso: en el primero corrige el nombre; en el segundo, reintenta más tarde.

\begin{lstlisting}
public interface PaisRepository {
    Pais buscar(String nombre) throws PaisNoEncontrado, PaisNoDisponible;
}
\end{lstlisting}

\texttt{CountriesDevPaisRepository} lee el JSON con Gson en un registro privado que copia la forma de la respuesta, y después lo convierte en \texttt{Pais}. Así, la forma del JSON de countries.dev no sale de esa clase. Tres detalles: el nombre se codifica en la URL (\texttt{Costa Rica} viaja como \texttt{Costa\%20Rica}); la búsqueda es por coincidencia parcial y \texttt{india} trae también el Territorio Británico del Océano Índico, así que gana el país cuyo nombre coincide exactamente; y la Antártida no tiene capital, así que \texttt{capital} puede llegar en \texttt{null}.

\subsection*{(b) El programa}

\texttt{App.main} recibe el nombre del país en los argumentos. Se ejecuta con:

\begin{lstlisting}
cd cliente-java
mvn -q compile exec:java -Dexec.args="Colombia"
\end{lstlisting}

Resultados del 5 de noviembre de 2026:

\begin{lstlisting}
$ Colombia
Colombia (capital Bogotá, 50.882.884 hab.): 30,0 °C, viento 1,2 km/h
$ Costa Rica
Costa Rica (capital San José, 5.094.114 hab.): 19,4 °C, viento 2,4 km/h
$ india
India (capital New Delhi, 1.380.004.385 hab.): 23,8 °C, viento 5,1 km/h
$ Antarctica
Antarctica (sin capital, 1.000 hab.): -34,0 °C, viento 18,9 km/h
\end{lstlisting}

Los números salen con el formato de Colombia (\texttt{Locale es-CO}): punto para los miles y coma para los decimales. Los 30~°C de Colombia no son los de Bogotá: \texttt{latlng} es el centro geográfico del país (\texttt{4, -72}, en los Llanos), no la capital. Con las coordenadas de Bogotá (\texttt{4.61, -74.08}), Open-Meteo daba 17,7~°C a la misma hora. La población es la de 2020, la que trae countries.dev; el ejemplo de la guía usaba un dato más reciente.

\subsection*{(c) Un país inexistente}

REST Countries v3.1 respondía \texttt{404 Not Found} a \texttt{/name/wakanda}. Hoy responde lo mismo que a cualquier otro país: \texttt{301} y, al seguir la redirección, \texttt{200} con el aviso de versión descontinuada. countries.dev responde \texttt{404 Not Found}, con un cuerpo que no es JSON: \texttt{Content-Type: text/plain} y el texto \texttt{Country not found}. Por eso el repositorio revisa el código antes de leer el cuerpo: si intentara convertir ese texto con Gson, fallaría. En Postman, las pruebas de esta petición verifican las tres cosas: el \texttt{404}, el \texttt{Content-Type} y el texto. Están en \texttt{postman/scripts-ejercicio2.js}, junto con las de las otras peticiones de la carpeta.

\captura{ej2-c-wakanda.png}{Figura 23. \texttt{GET https://countries.dev/name/wakanda}: 404 en texto plano.}

El \texttt{404} se convierte en \texttt{PaisNoEncontrado}, y \texttt{main} la atrapa y muestra un mensaje. Las demás fallas también terminan en un mensaje: lo comprobé sin red y con un servidor que acepta la conexión pero nunca responde (el programa se rindió a los 5,8 segundos):

\begin{lstlisting}[basicstyle=\ttfamily\footnotesize]
$ wakanda
No existe un país llamado "wakanda". Escriba el nombre en inglés, por ejemplo Spain o Germany.
$ Colombia        (sin red)
No se pudo completar la consulta: No hubo conexión con countries.dev. Intente más tarde.
$ Colombia        (servidor que no responde)
No se pudo completar la consulta: countries.dev no respondió en 5 segundos. Intente más tarde.
\end{lstlisting}

Los mensajes van por la salida de errores y el programa termina con código 1, para que un script que lo llame sepa que falló. Si no recibe un nombre, explica cómo usarlo y termina con código 2.

\subsection*{(d) Por qué \texttt{main} depende de la interfaz}

El principio de inversión de dependencias dice que la lógica de alto nivel no debe depender de los detalles, sino de abstracciones, y que la abstracción la define quien la usa. Aquí, \texttt{ClimaRepository} vive en el paquete \texttt{dominio} y dice lo que el dominio necesita (``el clima en estas coordenadas''), sin mencionar HTTP, URLs ni JSON. \texttt{OpenMeteoClimaRepository} vive en \texttt{infraestructura} e implementa esa interfaz. La dependencia apunta de la infraestructura hacia el dominio, no al revés.

El cálculo del resumen está en el método \texttt{resumen}, que recibe un \texttt{PaisRepository} y un \texttt{ClimaRepository}. \texttt{main} es el único lugar que conoce las clases concretas: las crea en dos líneas y se las pasa. Este ejercicio lo puso a prueba sin buscarlo: cuando REST Countries dejó de servir, el cambio fue escribir \texttt{CountriesDevPaisRepository} y cambiar una línea de \texttt{main}. \texttt{resumen}, \texttt{Pais} y la interfaz no se tocaron. La misma separación permite probar \texttt{resumen} con repositorios falsos que devuelven datos fijos, sin red.

En DDD, un Repository es la abstracción con la que el dominio obtiene sus objetos como si estuvieran en una colección, sin saber de dónde salen. Su interfaz pertenece al dominio y su implementación a la infraestructura. Aquí los datos no salen de una base de datos sino de una API, pero el papel es el mismo. Además, cada repositorio traduce el modelo ajeno al propio: el JSON de countries.dev (\texttt{name}, \texttt{latlng} como lista) se convierte en \texttt{Pais} (\texttt{nombre}, \texttt{lat}, \texttt{lon}) y los errores de red en excepciones del dominio. Así, los cambios de una API externa se quedan en una sola clase.

\subsection*{(e) En una app Android}

\begin{description}
  \raggedright
  \item[El hilo.] Android no deja hacer peticiones de red en el hilo principal, que es el que dibuja la pantalla: lanza \texttt{NetworkOnMainThreadException}. La petición tiene que correr en un hilo de fondo, por ejemplo con una corrutina en \texttt{viewModelScope} y \texttt{Dispatchers.IO}. El resultado se publica de vuelta para que la vista se actualice en el hilo principal. \texttt{viewModelScope} cancela la petición si el usuario cierra la pantalla.
  \item[El permiso.] En el \texttt{AndroidManifest.xml} hace falta \texttt{<uses-permission android:name="android.permission.INTERNET" />}. Es un permiso normal: se concede al instalar, sin preguntarle al usuario. Las dos APIs son HTTPS, y eso importa, porque desde Android 9 el tráfico HTTP sin cifrar está bloqueado por defecto.
  \item[El cliente HTTP.] Android no incluye \texttt{java.net.http.HttpClient}. Habría que cambiarlo por OkHttp o Retrofit, y ese cambio también queda dentro de las dos clases de infraestructura.
\end{description}

En MVVM, cada clase iría en una capa:

\begin{center}
\begin{tabular}{@{}>{\raggedright\arraybackslash}p{3.3cm}>{\raggedright\arraybackslash}p{11.6cm}@{}}
\hline
Capa & Qué hay \\
\hline
Model: dominio & \texttt{Clima}, \texttt{Pais}, \texttt{ClimaRepository}, \texttt{PaisRepository} y las excepciones. No cambian. \\
Model: datos & \texttt{OpenMeteoClimaRepository} y \texttt{CountriesDevPaisRepository}, reescritas con Retrofit u OkHttp. \\
ViewModel & Un \texttt{ClimaPaisViewModel} que recibe los repositorios por constructor, hace lo que hoy hace \texttt{resumen} y expone el estado de la pantalla (cargando, resultado o mensaje de error) como \texttt{StateFlow} o \texttt{LiveData}. Las excepciones se convierten aquí en mensajes, como hoy en \texttt{main}. \\
View & La \texttt{Activity}, el \texttt{Fragment} o la pantalla de Compose: un campo para el país y un botón. Observa el estado y lo dibuja; no hace peticiones. \\
\hline
\end{tabular}
\end{center}

\texttt{main} desaparece. La creación de los repositorios pasa a una \texttt{ViewModelProvider.Factory} o a un inyector de dependencias como Hilt.

\section*{Ejercicio 3. Los seis protocolos en la práctica}

\section*{Ejercicio 4. Diseño y documentación de una API REST}

\end{document}
